\documentclass[11pt,letterpaper]{article}
\usepackage[margin=1in]{geometry}
\usepackage{hyperref}
\usepackage{xcolor}
\usepackage{titlesec}
\usepackage{enumitem}
\usepackage{fancyhdr}
\usepackage{tocloft}
\usepackage{parskip}

% Colors
\definecolor{garnet}{RGB}{115,0,10}
\definecolor{darkgray}{RGB}{51,51,51}

% Hyperlink styling
\hypersetup{
    colorlinks=true,
    linkcolor=garnet,
    urlcolor=blue,
    citecolor=garnet
}

% Section formatting
\titleformat{\section}{\Large\bfseries\color{garnet}}{\thesection}{1em}{}
\titleformat{\subsection}{\large\bfseries\color{darkgray}}{\thesubsection}{1em}{}
\titleformat{\subsubsection}{\normalsize\bfseries}{\thesubsubsection}{1em}{}

% Header/Footer
\pagestyle{fancy}
\fancyhf{}
\fancyhead[L]{\textit{The Daily Gamecock Monthly}}
\fancyhead[R]{\textit{January 2026 Ideas}}
\fancyfoot[C]{\thepage}
\renewcommand{\headrulewidth}{0.4pt}

\begin{document}

% Title Page
\begin{titlepage}
    \centering
    \vspace*{2cm}
    {\Huge\bfseries\color{garnet} The Daily Gamecock\\Monthly Magazine\\[0.5cm]}
    {\LARGE Long-Form Opinion Piece Ideas\\[1cm]}
    {\Large January 2026 Edition\\[2cm]}

    \vfill

    {\large\textbf{10 Categories | 38+ Article Ideas}\\[0.5cm]}
    {\normalsize For 1,200--1,400 word evergreen analysis pieces\\[2cm]}

    \begin{tabular}{ll}
    \textbf{Categories:} & Student Debt \& Higher Ed Finance\\
    & AI \& Technology Transformation\\
    & Climate \& Environment\\
    & Campus Free Speech \& Academic Freedom\\
    & Mental Health \& Gen Z Culture\\
    & Economic Uncertainty \& Career Outcomes\\
    & Politics \& 2026 Midterms\\
    & Healthcare \& Public Health\\
    & Global Affairs \& American Isolationism\\
    & South Carolina \& USC Local Issues\\
    \end{tabular}

    \vfill

    {\small Compiled: January 14, 2026}
\end{titlepage}

\tableofcontents
\newpage

%==============================================================================
% CATEGORY 1: STUDENT DEBT & HIGHER ED FINANCE
%==============================================================================
\section{Student Debt \& Higher Education Finance}

\textit{Context: July 2026 brings major changes---SAVE plan ending, new borrowing caps, Grad Plus loans ending. 13 million could be in default by end of 2026. NIH/NSF funding cuts affecting research programs.}

%------------------------------------------------------------------------------
\newpage
\subsection{The Death of the SAVE Plan Will Push a Generation of USC Students Into Financial Ruin}

\subsubsection*{Core Argument}
The termination of the SAVE plan represents a deliberate policy choice to transfer the burden of higher education costs entirely onto students and families. USC students---particularly those from South Carolina's rural and low-income communities---will bear disproportionate consequences. This is not merely a budgetary adjustment; it is an abandonment of the social contract that once made American higher education accessible.

\subsubsection*{Angles to Explore}
\begin{enumerate}[noitemsep]
    \item \textbf{The Math of Devastation}: Under SAVE, payments were $\sim$\$150/month for median SC graduate; without it, \$350-400/month---a 130\% increase
    \item \textbf{South Carolina's Unique Vulnerability}: SC ranks 45th nationally in median household income; wages lag 15\% behind national average
    \item \textbf{The Equity Catastrophe}: First-generation students, Black students face compounded disadvantages
    \item \textbf{The Political Calculus}: Examine the ideological framework driving SAVE's termination
\end{enumerate}

\subsubsection*{Key Data Points}
\begin{itemize}[noitemsep]
    \item 13 million borrowers projected to be in default by end of 2026
    \item USC average graduate debt: $\sim$\$32,000 (undergrad), \$55,000+ (graduate)
    \item South Carolina median household income: \$59,318
    \item SAVE plan covered 8 million+ borrowers before termination
    \item USC Pell Grant recipient rate: $\sim$25\% of undergraduates
\end{itemize}

\subsubsection*{Sources to Interview}
Dr. Sandy Baum (Urban Institute), Dr. Fenaba Addo (UNC), USC Financial Aid Office, SC Student Loan Corporation, Student Borrower Protection Center, Ben Miller (former Biden administration)

\subsubsection*{Key URLs}
\begin{itemize}[noitemsep]
    \item \url{https://studentaid.gov/announcements-events/save-plan}
    \item \url{https://www.urban.org/policy-centers/cross-center-initiatives/building-americas-workforce}
    \item \url{https://protectborrowers.org/}
    \item \url{https://www.brookings.edu/articles/black-white-disparity-in-student-loan-debt-more-than-triples-after-graduation/}
\end{itemize}

%------------------------------------------------------------------------------
\newpage
\subsection{The Grad PLUS Loan Apocalypse Will Hollow Out USC's Research Mission}

\subsubsection*{Core Argument}
The elimination of Grad PLUS loans---combined with NIH and NSF funding cuts---will effectively end the accessibility of graduate education at research universities like USC. This isn't austerity; it's the intentional demolition of America's research infrastructure. USC's doctoral programs will become finishing schools for the wealthy while brilliant students from modest backgrounds are locked out.

\subsubsection*{Angles to Explore}
\begin{enumerate}[noitemsep]
    \item \textbf{The Hidden Economy of Graduate Education}: PhD stipends vs. cost of living in Columbia ($\sim$\$1,200+/month rent)
    \item \textbf{USC's Research Enterprise at Risk}: R1 status, Carnegie classification threatened
    \item \textbf{Disciplines That Will Disappear}: Social work, public health, education, humanities
    \item \textbf{The Wealth Filter}: Graduate education becomes hereditary privilege
\end{enumerate}

\subsubsection*{Key Data Points}
\begin{itemize}[noitemsep]
    \item NIH budget cuts: 20\%+ reduction affecting 2026-2027 grant cycles
    \item NSF funding cuts: $\sim$\$1.2 billion reduction
    \item USC awards $\sim$800 doctoral degrees annually
    \item Graduate student stipends at USC: \$20,000-\$28,000/year
    \item PhD completion rates: 50\% nationally
\end{itemize}

\subsubsection*{Sources to Interview}
USC Graduate Student Association, Dr. Hasan Jameel (USC VP Research), Department chairs in Biology, Chemistry, Social Work, Education; Dr. Colleen Campbell (Century Foundation), Council of Graduate Schools

\subsubsection*{Key URLs}
\begin{itemize}[noitemsep]
    \item \url{https://www.science.org/news}
    \item \url{https://cgsnet.org/}
    \item \url{https://cew.georgetown.edu/}
    \item \url{https://www.sc.edu/study/colleges_schools/graduate_school/}
\end{itemize}

%------------------------------------------------------------------------------
\newpage
\subsection{South Carolina's Refusal to Expand State Aid Will Make USC Unaffordable for the Students It Was Built to Serve}

\subsubsection*{Core Argument}
While federal student aid collapses, South Carolina has a choice: expand state grants and need-based aid, or watch USC become inaccessible to middle-class and low-income South Carolinians. The state's current trajectory---minimal need-based aid, heavy reliance on merit scholarships that disproportionately benefit wealthier families---suggests it will choose abandonment.

\subsubsection*{Angles to Explore}
\begin{enumerate}[noitemsep]
    \item \textbf{SC's Higher Ed Funding Crisis}: State appropriations cover only $\sim$10\% of USC's budget (down from 50\%+ in 1980s)
    \item \textbf{The LIFE Scholarship Scam}: Merit-based aid disproportionately benefits wealthier families
    \item \textbf{What Other States Are Doing}: Georgia HOPE, New Mexico tuition-free, NC need-based aid
    \item \textbf{The Legislative Path Forward}: Specific policy proposals with cost estimates
\end{enumerate}

\subsubsection*{Key Data Points}
\begin{itemize}[noitemsep]
    \item SC state funding per FTE student: $\sim$\$6,500 (national average: \$9,000+)
    \item LIFE Scholarship: \$5,000/year (hasn't increased significantly while tuition rose 40\%+)
    \item SC need-based grant funding: under \$30 million annually (vs. NC's \$200+ million)
    \item USC in-state tuition: \$13,000+/year (up from \$5,000 in 2000)
\end{itemize}

\subsubsection*{Sources to Interview}
SC Commission on Higher Education, state legislators on education committees, Dr. Sara Goldrick-Rab (Temple), TICAS, USC Office of the President, SC Chamber of Commerce

\subsubsection*{Key URLs}
\begin{itemize}[noitemsep]
    \item \url{https://www.che.sc.gov/}
    \item \url{https://shef.sheeo.org/}
    \item \url{https://ticas.org/}
    \item \url{https://hope.temple.edu/}
\end{itemize}

%------------------------------------------------------------------------------
\newpage
\subsection{USC Students Should Collectively Refuse to Pay Until the System Changes}

\subsubsection*{Core Argument}
Individual financial planning advice cannot solve a systemic crisis. With 13 million borrowers heading toward default, the student debt system is already collapsing; students should accelerate that collapse through collective action. A debt strike and coordinated non-payment are not irresponsible---they are the only rational response to an irrational system.

\subsubsection*{Angles to Explore}
\begin{enumerate}[noitemsep]
    \item \textbf{The Precedent}: Corinthian Fifteen won billions in debt cancellation through collective refusal
    \item \textbf{The Moral Case for Non-Payment}: Contracts obtained through misrepresentation are voidable
    \item \textbf{What USC Students Can Actually Do}: Form Debt Collective chapter, coordinate with SG across SC
    \item \textbf{The Counterarguments and Why They're Wrong}: Address credit, responsibility objections
\end{enumerate}

\subsubsection*{Key Data Points}
\begin{itemize}[noitemsep]
    \item Total US student debt: \$1.77 trillion
    \item Number of borrowers: 43 million+
    \item Debt cancellation won by Debt Collective organizing: \$2+ billion
    \item 34\% of borrowers would join debt strike if coordinated (polling data)
\end{itemize}

\subsubsection*{Sources to Interview}
The Debt Collective (Astra Taylor, Thomas Gokey), Dr. Hannah Appel (UCLA), Strike Debt activists, legal experts on debt collection, Consumer Financial Protection Bureau

\subsubsection*{Key URLs}
\begin{itemize}[noitemsep]
    \item \url{https://debtcollective.org/}
    \item \url{https://strike.debtcollective.org/}
    \item \url{https://www.consumerfinance.gov/data-research/student-lending/}
    \item \url{https://www.nclc.org/issues/student-loan-law/}
\end{itemize}

%==============================================================================
% CATEGORY 2: AI & TECHNOLOGY TRANSFORMATION
%==============================================================================
\newpage
\section{AI \& Technology Transformation}

\textit{Context: 92\% of students now use AI. Purdue requires AI literacy for graduation. 70\% of students worry AI will eliminate entry-level jobs. AI wage premium at 56\%.}

%------------------------------------------------------------------------------
\newpage
\subsection{The AI Literacy Requirement Purdue Got Right That USC Cannot Afford to Ignore}

\subsubsection*{Core Argument}
South Carolina's flagship university is failing its students by not implementing mandatory AI literacy requirements. While Purdue has made AI competency a graduation requirement, USC continues to treat AI education as optional, leaving graduates unprepared for a job market that now pays a 56\% wage premium for AI skills.

\subsubsection*{Angles to Explore}
\begin{enumerate}[noitemsep]
    \item \textbf{The Economic Imperative}: 56\% wage premium represents the difference between middle-class jobs and being locked out
    \item \textbf{The Equity Dimension}: Leaving AI literacy to individual initiative creates two-tiered system
    \item \textbf{The Departmental Resistance Problem}: Why USC hasn't followed Purdue's lead
    \item \textbf{What ``AI Literacy'' Actually Means}: Beyond ChatGPT familiarity to genuine competency
\end{enumerate}

\subsubsection*{Key Data Points}
\begin{itemize}[noitemsep]
    \item 56\% wage premium for AI-skilled workers
    \item 92\% of students currently use AI tools
    \item Purdue's graduation requirement implementation (2024-2025)
    \item 70\% of students worry AI will eliminate entry-level jobs
\end{itemize}

\subsubsection*{Sources to Interview}
Dr. Amit Sheth (USC AI Institute), Dr. Forest Agostinelli (USC CS), Purdue Provost's office, USC Career Center, Dr. Ethan Mollick (Wharton), local SC employers (Boeing, Prisma Health, Dominion Energy)

\subsubsection*{Key URLs}
\begin{itemize}[noitemsep]
    \item \url{https://www.purdue.edu/provost/}
    \item \url{https://www.weforum.org/publications/the-future-of-jobs-report-2023/}
    \item \url{https://aiisc.ai/}
    \item \url{https://www.oneusefulthing.org/}
\end{itemize}

%------------------------------------------------------------------------------
\newpage
\subsection{AI Detection Tools Are the New Racial Profiling and USC Must Abandon Them}

\subsubsection*{Core Argument}
AI detection tools are demonstrably biased against non-native English speakers, falsely flagging international students at significantly higher rates. USC's continued reliance on these tools constitutes technological discrimination that undermines the university's commitment to diversity and inclusion.

\subsubsection*{Angles to Explore}
\begin{enumerate}[noitemsep]
    \item \textbf{The Technical Failure}: 60\%+ false positive rates for non-native English writing
    \item \textbf{The Human Cost at USC}: Psychological toll, power imbalance for 1,800+ international students
    \item \textbf{The Legal and Ethical Exposure}: Potential Title VI liability from disparate impact
    \item \textbf{Better Alternatives Exist}: Process-based assessment, oral examinations, assignment redesign
\end{enumerate}

\subsubsection*{Key Data Points}
\begin{itemize}[noitemsep]
    \item Studies showing 60\%+ false positive rates for non-native speakers
    \item USC international student enrollment: $\sim$1,800+
    \item Turnitin's ``98\% confidence'' means 2\% false positives
    \item GPTZero and detector accuracy studies from Stanford, University of Maryland
\end{itemize}

\subsubsection*{Sources to Interview}
USC International Student Services, Dr. Deborah Weber-Wulff (HTW Berlin), USC Writing Center, Dr. Emily M. Bender (UW), legal experts on Title VI

\subsubsection*{Key URLs}
\begin{itemize}[noitemsep]
    \item \url{https://hai.stanford.edu/}
    \item \url{https://www.turnitin.com/products/features/ai-writing-detection}
    \item \url{https://gptzero.me/}
    \item \url{https://www.sc.edu/about/offices_and_divisions/student_conduct_and_academic_integrity/}
\end{itemize}

%------------------------------------------------------------------------------
\newpage
\subsection{Your Major Won't Protect You From AI and Career Centers Are Lying By Omission}

\subsubsection*{Core Argument}
USC's career counseling apparatus is failing students by continuing to operate on pre-AI assumptions about ``safe'' career paths. The comfortable advice that certain majors are ``AI-proof'' is dangerously outdated. Career preparation in 2026 requires a fundamentally different approach.

\subsubsection*{Angles to Explore}
\begin{enumerate}[noitemsep]
    \item \textbf{The Myth of the ``AI-Proof'' Major}: Pre-law, healthcare, engineering, creative fields all face disruption
    \item \textbf{The Entry-Level Collapse}: 70\% worried about entry-level elimination are right
    \item \textbf{What Honest Career Counseling Looks Like}: Expect to change every 2-3 years, build complementary skills
    \item \textbf{The USC Employer Disconnect}: What employers actually want vs. student preparation
\end{enumerate}

\subsubsection*{Key Data Points}
\begin{itemize}[noitemsep]
    \item 70\% of students worry AI will eliminate entry-level jobs
    \item 56\% AI wage premium
    \item Job posting analysis showing AI-skill requirements in ``non-technical'' fields
    \item USC Career Center programming analysis
\end{itemize}

\subsubsection*{Sources to Interview}
USC Career Center director, USC College of Engineering career services, Moore School career placement, employers recruiting at USC, Dr. Carl Benedikt Frey (Oxford), recent USC graduates

\subsubsection*{Key URLs}
\begin{itemize}[noitemsep]
    \item \url{https://www.bls.gov/}
    \item \url{https://economicgraph.linkedin.com/}
    \item \url{https://www.oxfordmartin.ox.ac.uk/future-of-work/}
    \item \url{https://sc.edu/career}
\end{itemize}

%------------------------------------------------------------------------------
\newpage
\subsection{The 2026 Quantum Computing Milestone Changes Everything and USC Students Are Woefully Unprepared}

\subsubsection*{Core Argument}
The quantum computing milestone achieved in 2026 isn't just a physics breakthrough---it's the starting gun for transformation that will make AI disruption look like a warm-up. USC has perhaps 5-10 years before quantum computing transforms multiple industries, yet offers virtually no quantum literacy preparation outside physics and CS.

\subsubsection*{Angles to Explore}
\begin{enumerate}[noitemsep]
    \item \textbf{What the 2026 Milestone Means}: Quantum advantage, practical applications timeline
    \item \textbf{Industries About to Be Disrupted}: Cybersecurity, pharma, finance, logistics
    \item \textbf{The Security Emergency}: ``Harvest now, decrypt later'' attacks on encrypted data
    \item \textbf{Why USC Should Lead}: Strong physics program could become quantum education leader
\end{enumerate}

\subsubsection*{Key Data Points}
\begin{itemize}[noitemsep]
    \item 2026 quantum computing milestone specifics
    \item Quantum computing market: \$65+ billion by 2030
    \item Current USC quantum-related course offerings
    \item NIST post-quantum cryptography standards timeline
\end{itemize}

\subsubsection*{Sources to Interview}
USC Physics Department quantum researchers, USC CS faculty on quantum algorithms, Dr. John Preskill (Caltech), IBM Quantum, NIST post-quantum cryptography team

\subsubsection*{Key URLs}
\begin{itemize}[noitemsep]
    \item \url{https://www.quantum.gov/}
    \item \url{https://www.ibm.com/quantum}
    \item \url{https://csrc.nist.gov/projects/post-quantum-cryptography}
    \item \url{https://www.technologyreview.com/topic/computing/quantum-computing/}
\end{itemize}

%==============================================================================
% CATEGORY 3: CLIMATE & ENVIRONMENT
%==============================================================================
\newpage
\section{Climate \& Environment}

\textit{Context: US withdrew from Paris Agreement, UNFCCC, IPCC. US fell to 65th in Climate Performance Index. 2025 second-hottest year ever. South Carolina coastal vulnerability.}

%------------------------------------------------------------------------------
\newpage
\subsection{South Carolina's Coastline Will Disappear Within Our Lifetimes Unless We Act Now}

\subsubsection*{Core Argument}
The federal government's withdrawal from international climate agreements has effectively abandoned South Carolina's 187-mile coastline to accelerating sea-level rise. USC students who grew up visiting Myrtle Beach, Hilton Head, and Charleston will watch these places become uninhabitable unless state and local action fills the void.

\subsubsection*{Angles to Explore}
\begin{enumerate}[noitemsep]
    \item \textbf{The Economics of Coastal Collapse}: SC's \$24 billion coastal tourism economy, insurance market failures
    \item \textbf{Environmental Justice}: Gullah-Geechee communities face displacement without resources
    \item \textbf{The Infrastructure Reckoning}: Charleston 89+ days of tidal flooding annually (up from 2 in 1970s)
    \item \textbf{Generational Wealth Transfer in Reverse}: Inheriting underwater mortgages, not valuable property
\end{enumerate}

\subsubsection*{Key Data Points}
\begin{itemize}[noitemsep]
    \item Sea levels along SC coast risen 8-10 inches since 1950
    \item Charleston projects 2+ feet additional rise by 2050
    \item Hurricane damage costs in SC: \$25+ billion (2015-2024)
    \item 540,000+ SC properties face substantial flood risk
    \item Coastal tourism generates 113,000 jobs
\end{itemize}

\subsubsection*{Sources to Interview}
Dr. Scott Curtis (USC Hazards Institute), Dr. Marcia Wilson (USC Geography), Jason Crowley (Coastal Conservation League), Mayor Tecklenburg (Charleston), Gullah/Geechee Heritage Corridor Commission

\subsubsection*{Key URLs}
\begin{itemize}[noitemsep]
    \item \url{https://coast.noaa.gov/slr/}
    \item \url{https://floodfactor.com/}
    \item \url{https://www.scseagrant.org/}
    \item \url{https://sealevel.climatecentral.org/}
\end{itemize}

%------------------------------------------------------------------------------
\newpage
\subsection{The Green Economy Is Coming Whether Washington Likes It or Not}

\subsubsection*{Core Argument}
The federal government's rescission of clean energy tax credits represents a strategic miscalculation. However, the green economy transition is driven by fundamental economics, not policy preference. USC students who position themselves now will thrive while those waiting for fossil fuel nostalgia will be left behind.

\subsubsection*{Angles to Explore}
\begin{enumerate}[noitemsep]
    \item \textbf{The Inevitability Argument}: Solar/wind now cheapest, battery costs down 97\% since 1991
    \item \textbf{SC's Manufacturing Opportunity}: BMW Spartanburg, Volvo Charleston transitioning to EVs
    \item \textbf{The Skills Gap as Opportunity}: Not enough solar installers, wind technicians, EV mechanics
    \item \textbf{Entrepreneurship and Innovation}: Federal retreat creates market opportunities
\end{enumerate}

\subsubsection*{Key Data Points}
\begin{itemize}[noitemsep]
    \item Global renewable investment: \$1.8 trillion in 2024
    \item Solar jobs grew 5.3\% in 2024 despite policy uncertainty
    \item SC ranks 15th nationally in solar potential
    \item Clean energy jobs pay 25\% more than national median
    \item SC manufacturing employs 250,000+ workers
\end{itemize}

\subsubsection*{Sources to Interview}
USC SmartState Center researchers, Moore School sustainability initiatives, Hamilton Davis (SC Solar Business Alliance), BMW/Volvo representatives, Dominion Energy SC

\subsubsection*{Key URLs}
\begin{itemize}[noitemsep]
    \item \url{https://about.bnef.com/}
    \item \url{https://www.bls.gov/green/}
    \item \url{https://www.irena.org/}
    \item \url{https://www.scsolaralliance.org/}
\end{itemize}

%------------------------------------------------------------------------------
\newpage
\subsection{USC's Sustainability Efforts Are Performative Without Institutional Divestment}

\subsubsection*{Core Argument}
USC promotes sustainability initiatives while investing endowment funds in fossil fuel companies whose business model requires climate destruction. This hypocrisy undermines every solar panel installed and every sustainability course offered.

\subsubsection*{Angles to Explore}
\begin{enumerate}[noitemsep]
    \item \textbf{Divestment Movement's Success}: 1,500+ institutions (\$40+ trillion) committed
    \item \textbf{Quantifying USC's Climate Impact}: Endowment investments vs. campus emissions reductions
    \item \textbf{Student and Faculty Organizing}: History of divestment advocacy at USC
    \item \textbf{The ``Engagement'' Excuse}: Shareholder advocacy has failed for decades
\end{enumerate}

\subsubsection*{Key Data Points}
\begin{itemize}[noitemsep]
    \item USC endowment: $\sim$\$1.1 billion
    \item 1,500+ institutions globally committed to divestment
    \item UC divested \$126 billion, reported equal or better returns
    \item Carbon Tracker: 80\% of fossil reserves must stay unburned
    \item USC carbon neutrality goal: 2045
\end{itemize}

\subsubsection*{Sources to Interview}
USC Foundation Board, USC President Amiridis, Bill Cain (USC Sustainability), Ellen Dorsey (Wallace Global Fund), Bill McKibben (350.org), representatives from divested institutions

\subsubsection*{Key URLs}
\begin{itemize}[noitemsep]
    \item \url{https://divestmentdatabase.org/}
    \item \url{https://350.org/campaigns/divestment/}
    \item \url{https://carbontracker.org/}
    \item \url{https://stars.aashe.org/}
\end{itemize}

%------------------------------------------------------------------------------
\newpage
\subsection{Climate Denial in the State Legislature Is a Direct Threat to Every USC Student's Future}

\subsubsection*{Core Argument}
South Carolina's legislature has consistently blocked climate action and prioritized fossil fuel interests. As federal climate policy collapses, state-level action becomes the only meaningful pathway---and SC's legislative environment is actively hostile to that pathway.

\subsubsection*{Angles to Explore}
\begin{enumerate}[noitemsep]
    \item \textbf{Follow the Money}: Fossil fuel contributions to SC legislators who block climate action
    \item \textbf{V.C. Summer Debacle}: \$9 billion ratepayer loss from failed nuclear project
    \item \textbf{What Other States Are Doing}: NC, GA, VA clean energy policies
    \item \textbf{Youth Political Power}: USC's 35,000+ students as decisive political force
\end{enumerate}

\subsubsection*{Key Data Points}
\begin{itemize}[noitemsep]
    \item SC ranks 44th in solar policy
    \item V.C. Summer cost ratepayers \$9 billion
    \item SC has no renewable portfolio standards
    \item USC student body: 35,000+ potential voters in Richland County
\end{itemize}

\subsubsection*{Sources to Interview}
SC Conservation Voters, Energy and Policy Institute, Dr. Robert Oldendick (USC Political Science), USC College Democrats/Republicans, state legislators, Climate Cabinet

\subsubsection*{Key URLs}
\begin{itemize}[noitemsep]
    \item \url{https://ethics.sc.gov/}
    \item \url{https://www.scstatehouse.gov/}
    \item \url{https://www.energyandpolicy.org/}
    \item \url{https://www.climatecabinet.org/}
\end{itemize}

%==============================================================================
% CATEGORY 4: CAMPUS FREE SPEECH & ACADEMIC FREEDOM
%==============================================================================
\newpage
\section{Campus Free Speech \& Academic Freedom}

\textit{Context: FIRE recorded 525 Scholar Under Fire incidents. Government censorship up to 1/3 of cases. SC politicians calling to eliminate tenure. DEI dismantled at 437 campuses.}

%------------------------------------------------------------------------------
\newpage
\subsection{South Carolina's War on Tenure Will Destroy the University System It Claims to Protect}

\subsubsection*{Core Argument}
SC legislators' push to eliminate tenure represents a fundamental misunderstanding of how academic excellence is produced. Their efforts will drive top faculty to other states, diminish USC's research output, and harm economic development goals. The attack on tenure is a solution in search of a problem.

\subsubsection*{Angles to Explore}
\begin{enumerate}[noitemsep]
    \item \textbf{Economic Self-Sabotage}: Faculty recruitment, research grants, startup partnerships at risk
    \item \textbf{The ``Accountability'' Myth}: Post-tenure review already exists
    \item \textbf{Chilling Effect on Controversial Research}: Climate, public health, criminal justice research threatened
    \item \textbf{National Competition for Talent}: Faculty being recruited away
\end{enumerate}

\subsubsection*{Key Data Points}
\begin{itemize}[noitemsep]
    \item FIRE: 525 ``Scholar Under Fire'' incidents
    \item AAUP data on faculty salaries and tenure rates by state
    \item USC faculty retention statistics
    \item SC ranks near bottom in higher education funding per student
\end{itemize}

\subsubsection*{Sources to Interview}
Dr. Irene Mulvey (AAUP President), USC Faculty Senate, Dr. Hans-Joerg Tiede (AAUP Research), SC legislators, faculty at Wisconsin/UNC who experienced tenure weakening

\subsubsection*{Key URLs}
\begin{itemize}[noitemsep]
    \item \url{https://www.thefire.org/research-learn/scholars-under-fire}
    \item \url{https://www.aaup.org/issues/tenure}
    \item \url{https://www.che.sc.gov/}
    \item \url{https://sc.edu/about/offices_and_divisions/provost/docs/faculty_manual.pdf}
\end{itemize}

%------------------------------------------------------------------------------
\newpage
\subsection{The Harvard Deportation Threat Should Terrify Every International Student at USC}

\subsubsection*{Core Argument}
The federal government's visa revocation based on social media posts---not criminal conduct---represents a fundamental shift for international students. USC's 2,500+ international students now face an impossible choice between authentic intellectual engagement and self-protective silence.

\subsubsection*{Angles to Explore}
\begin{enumerate}[noitemsep]
    \item \textbf{The New Rules of Engagement}: What activities now carry deportation risk?
    \item \textbf{USC's International Student Community Response}: Self-censorship in real-time
    \item \textbf{Institutional Responsibility}: What has USC communicated to international students?
    \item \textbf{The Two-Tiered Campus}: Different speech rights for citizens vs. non-citizens
\end{enumerate}

\subsubsection*{Key Data Points}
\begin{itemize}[noitemsep]
    \item USC international students: 2,500+
    \item Percentage of USC graduate programs relying on international enrollment
    \item Economic contribution of international students to SC
    \item Visa revocations nationally in 2024-2025
\end{itemize}

\subsubsection*{Sources to Interview}
USC International Student Services, immigration attorneys, USC international students (confidentially), ACLU Immigrants' Rights Project, Dr. Sarah McLaughlin (FIRE)

\subsubsection*{Key URLs}
\begin{itemize}[noitemsep]
    \item \url{https://www.nafsa.org/}
    \item \url{https://www.thefire.org/}
    \item \url{https://www.aclu.org/know-your-rights/immigrants-rights}
    \item \url{https://www.ice.gov/sevis}
\end{itemize}

%------------------------------------------------------------------------------
\newpage
\subsection{Dismantling DEI Hasn't Made Campuses Freer---It's Just Shifted Who Gets Silenced}

\subsubsection*{Core Argument}
The nationwide dismantling of DEI programs---framed as a victory for free speech---has not actually expanded acceptable discourse. It has merely changed which perspectives face institutional hostility. The ``anti-DEI'' movement has become its own orthodoxy.

\subsubsection*{Angles to Explore}
\begin{enumerate}[noitemsep]
    \item \textbf{The Censorship Shell Game}: ``Anti-DEI'' policies restricting speech about racism
    \item \textbf{SC Legislative Context}: How anti-DEI legislation implemented at USC
    \item \textbf{Student Experience Gap}: Effects on minority, LGBTQ+, first-gen students
    \item \textbf{Intellectual Dishonesty of ``Neutrality''}: Absence of programming is itself a position
\end{enumerate}

\subsubsection*{Key Data Points}
\begin{itemize}[noitemsep]
    \item 437+ campuses dismantling DEI (Chronicle of Higher Education)
    \item 135 anti-DEI bills introduced; 29 signed into law
    \item USC DEI program funding before and after changes
    \item FIRE campus climate survey data
\end{itemize}

\subsubsection*{Sources to Interview}
USC faculty in affected departments, Ilya Shapiro (Manhattan Institute), Dr. Shaun Harper (USC Race and Equity Center), students from diverse backgrounds, FIRE staff

\subsubsection*{Key URLs}
\begin{itemize}[noitemsep]
    \item \url{https://www.chronicle.com/article/tracking-higher-eds-dismantling-of-dei}
    \item \url{https://www.thefire.org/college-free-speech-rankings}
    \item \url{https://knightfoundation.org/topics/free-expression/}
    \item \url{https://www.scstatehouse.gov/}
\end{itemize}

%------------------------------------------------------------------------------
\newpage
\subsection{USC's Controversial Speaker Dilemma Reveals That Nobody Actually Believes in Free Speech}

\subsubsection*{Core Argument}
The debate over controversial speakers consistently conflates two questions: whether someone has the right to speak and whether a university must provide a platform. Students protesting are exercising free speech, not violating it. The real threat comes from administrative decisions that cave to pressure or fail to articulate coherent principles.

\subsubsection*{Angles to Explore}
\begin{enumerate}[noitemsep]
    \item \textbf{Platform vs. Censorship Distinction}: Censorship, deplatforming, and protest are different things
    \item \textbf{USC's 2025 Speaker Controversies}: Specific case studies
    \item \textbf{The ``Both Sides'' Trap}: Not all positions are legitimate academic perspectives
    \item \textbf{A Principled Framework}: Maximum space for controversial speech AND vigorous protest
\end{enumerate}

\subsubsection*{Key Data Points}
\begin{itemize}[noitemsep]
    \item Number of speaker ``disinvitations'' nationally (FIRE database)
    \item USC speaker events in 2025 and outcomes
    \item Security cost figures for controversial events
    \item Student survey data on attitudes toward controversial speakers
\end{itemize}

\subsubsection*{Sources to Interview}
USC student organization leaders, Suzanne Nossel (PEN America), FIRE representatives, USC administration, students who protested speakers, Greg Lukianoff (FIRE)

\subsubsection*{Key URLs}
\begin{itemize}[noitemsep]
    \item \url{https://www.thefire.org/research-learn/campus-disinvitation-database}
    \item \url{https://pen.org/issue/free-expression-and-education/}
    \item \url{https://heterodoxacademy.org/}
    \item \url{https://www.aclu.org/issues/free-speech/rights-protesters}
\end{itemize}

%==============================================================================
% CATEGORY 5: MENTAL HEALTH & GEN Z CULTURE
%==============================================================================
\newpage
\section{Mental Health \& Gen Z Culture}

\textit{Context: ``Social exit'' movement---abandoning social media. Only 27\% of undergrads rate mental health as above average. 68\% say mental health impacts academics. Counseling center utilization up 20-30\%.}

%------------------------------------------------------------------------------
\newpage
\subsection{The Quiet Rebellion of Logging Off Is the Most Radical Act of Our Generation}

\subsubsection*{Core Argument}
The ``social exit'' movement among Gen Z is not retreat---it is a deliberate, politically conscious rejection of the attention economy that has monetized their mental health since childhood. When USC students delete Instagram or switch to flip phones, they are engaging in protest.

\subsubsection*{Angles to Explore}
\begin{enumerate}[noitemsep]
    \item \textbf{Economics of Attention Theft}: Platforms extracting value from Gen Z's developmental years
    \item \textbf{The Privilege Problem}: Who can afford to log off?
    \item \textbf{SC's Legislative Response}: Social Media Regulation Act
    \item \textbf{Campus Culture Without the Feed}: Student orgs, Greek life when segment refuses digital engagement
\end{enumerate}

\subsubsection*{Key Data Points}
\begin{itemize}[noitemsep]
    \item Only 27\% of undergrads rate mental health as above average
    \item 68\% report mental health impacts academics
    \item USC Counseling Center utilization up 20-30\%
    \item Pew: 54\% of teens say hard to give up social media (down from 2018)
\end{itemize}

\subsubsection*{Sources to Interview}
Dr. Jean Twenge (SDSU), USC Counseling and Psychiatry Services, Center for Humane Technology, USC students who exited social media, Dr. Anna Lembke (Stanford)

\subsubsection*{Key URLs}
\begin{itemize}[noitemsep]
    \item \url{https://www.humanetech.com/}
    \item \url{https://www.apa.org/topics/social-media-internet/health-advisory-adolescent-social-media-use}
    \item \url{https://www.pewresearch.org/internet/}
    \item \url{https://www.jeantwenge.com/}
\end{itemize}

%------------------------------------------------------------------------------
\newpage
\subsection{We Moved Back Home and It Saved Us From the Loneliness Economy}

\subsubsection*{Core Argument}
The 31\% of young adults who moved back home are not failures of independence but pioneers of a necessary cultural correction. Gen Z's return to multigenerational living represents rejection of atomized, hyper-individualist lifestyle that produced an epidemic of loneliness.

\subsubsection*{Angles to Explore}
\begin{enumerate}[noitemsep]
    \item \textbf{The Mythology of Independence}: Challenge American narrative of geographic separation
    \item \textbf{Economics Behind the Exodus}: Housing costs make independence financially devastating
    \item \textbf{Loneliness as Public Health Crisis}: Surgeon General's 2023 advisory
    \item \textbf{What This Means for ``The College Experience''}: Should USC expand commuter support?
\end{enumerate}

\subsubsection*{Key Data Points}
\begin{itemize}[noitemsep]
    \item 31\% of young adults moved back home
    \item Loneliness mortality impact: equivalent to smoking 15 cigarettes daily
    \item USC housing costs vs. median family income in SC
    \item Pew: more 18-29 year olds with parents than since Great Depression
\end{itemize}

\subsubsection*{Sources to Interview}
U.S. Surgeon General's Office, USC Housing and Residential Life, Dr. Bella DePaulo, USC students living at home, NAMI South Carolina, Economic Policy Institute

\subsubsection*{Key URLs}
\begin{itemize}[noitemsep]
    \item \url{https://www.hhs.gov/surgeongeneral/priorities/connection/index.html}
    \item \url{https://www.pewresearch.org/}
    \item \url{https://namisc.org/}
    \item \url{https://housing.sc.edu/}
\end{itemize}

%------------------------------------------------------------------------------
\newpage
\subsection{The Campus Counseling Crisis Proves We Cannot Therapy Our Way Out of a Sick System}

\subsubsection*{Core Argument}
The 20-30\% increase in counseling utilization reflects a generation correctly identifying something is deeply wrong---but individual therapy cannot solve systemic problems. The framing of Gen Z's distress as clinical obscures structural causes: precarious economics, climate anxiety, political instability.

\subsubsection*{Angles to Explore}
\begin{enumerate}[noitemsep]
    \item \textbf{Limits of the Therapeutic Model}: Therapist can't solve \$40,000 debt load
    \item \textbf{SC's Youth Mental Health Act}: Addressing symptoms while ignoring causes?
    \item \textbf{Professionalization of Coping}: Wellness industry profiting from distress
    \item \textbf{Systemic Solutions}: What if USC invested equally in addressing causes?
\end{enumerate}

\subsubsection*{Key Data Points}
\begin{itemize}[noitemsep]
    \item USC Counseling Center utilization up 20-30\%
    \item 68\% report mental health impacts academics
    \item Counselor-to-student ratio at USC vs. recommended
    \item Average student debt for USC graduates
\end{itemize}

\subsubsection*{Sources to Interview}
USC Counseling leadership, Dr. Jonathan Shedler, USC Student Government, SC legislators (Youth Mental Health Act), Active Minds USC, Dr. James Davies (``Sedated'' author)

\subsubsection*{Key URLs}
\begin{itemize}[noitemsep]
    \item \url{https://ccmh.psu.edu/}
    \item \url{https://sc.edu/about/offices_and_divisions/student_health_services/medical-services/counseling-and-psychiatry/}
    \item \url{https://www.activeminds.org/}
    \item \url{https://www.naspa.org/}
\end{itemize}

%------------------------------------------------------------------------------
\newpage
\subsection{Knitting Circles and Flip Phones Will Not Save Us But They Point Toward What Might}

\subsubsection*{Core Argument}
The return to analog hobbies represents a genuine insight (tactile engagement supports mental health) wrapped in a troubling illusion (individual lifestyle changes can't substitute for collective action). USC students rediscovering analog life are diagnosing correctly while prescribing insufficiently.

\subsubsection*{Angles to Explore}
\begin{enumerate}[noitemsep]
    \item \textbf{Genuine Benefits of Analog}: Research supports tactile hobbies, flow states
    \item \textbf{The Consumption Trap}: Analog hobbies still require purchasing
    \item \textbf{Individual Solutions to Collective Problems}: 68\% can't all take up embroidery
    \item \textbf{Collective ``Analog'' Action}: What if knitting circle became union drive?
\end{enumerate}

\subsubsection*{Key Data Points}
\begin{itemize}[noitemsep]
    \item Screen time dropping after years of increase
    \item Vinyl sales, film photography, craft supplies market growth
    \item USC student organization data: growth in craft clubs vs. activist organizations
    \item Research on flow states and mental health (Csikszentmihalyi)
\end{itemize}

\subsubsection*{Sources to Interview}
USC craft/hobby organizations, Csikszentmihalyi research successors, Dr. Juliet Schor (Boston College), labor organizers, USC Recreation and Wellness

\subsubsection*{Key URLs}
\begin{itemize}[noitemsep]
    \item \url{https://www.pursuit-of-happiness.org/history-of-happiness/mihaly-csikszentmihalyi/}
    \item \url{https://www.bc.edu/bc-web/schools/morrissey/departments/sociology/people/faculty-directory/juliet-schor.html}
    \item \url{https://www.riaa.com/}
    \item \url{https://sc.edu/about/offices_and_divisions/leadership_and_service_center/student-organizations/}
\end{itemize}

%==============================================================================
% CATEGORY 6: ECONOMIC UNCERTAINTY & CAREER OUTCOMES
%==============================================================================
\newpage
\section{Economic Uncertainty \& Career Outcomes}

\textit{Context: 1.6\% hiring increase projected for Class of 2026. 41.8\% underemployment rate (highest since 2020). 22-27 year olds have higher unemployment than general workforce for first time.}

%------------------------------------------------------------------------------
\newpage
\subsection{The Diploma Deception Has Finally Caught Up With America's College Students}

\subsubsection*{Core Argument}
The implicit social contract that a four-year degree guarantees middle-class employment has fundamentally collapsed. The 41.8\% underemployment rate isn't a temporary blip; it's the new equilibrium that exposes how higher education failed to adapt. Students aren't failing the job market; institutions failed to prepare them.

\subsubsection*{Angles to Explore}
\begin{enumerate}[noitemsep]
    \item \textbf{The Credential Inflation Trap}: Degrees required even as skills needed remained static
    \item \textbf{USC's Major-Specific Vulnerability}: Which colleges face steepest cliffs
    \item \textbf{The Hidden Curriculum Gap}: Professional skills employers actually screen for
    \item \textbf{South Carolina's Paradox}: GDP growth in manufacturing, not bachelor's-required positions
\end{enumerate}

\subsubsection*{Key Data Points}
\begin{itemize}[noitemsep]
    \item 41.8\% underemployment rate (highest since 2020)
    \item 1.6\% projected hiring increase for Class of 2026
    \item 22-27 year olds: higher unemployment than general workforce (first time)
    \item Fed data: 45\% of underemployed remain so after five years
\end{itemize}

\subsubsection*{Sources to Interview}
Dr. Douglas Webber (Temple), USC Career Center, Burning Glass Institute, SC DEW economists, Dr. Anthony Carnevale (Georgetown), recent USC alumni working outside field

\subsubsection*{Key URLs}
\begin{itemize}[noitemsep]
    \item \url{https://www.newyorkfed.org/research/college-labor-market/college-labor-market_underemployment_rates}
    \item \url{https://www.burningglassinstitute.org/}
    \item \url{https://cew.georgetown.edu/}
    \item \url{https://www.naceweb.org/job-market/trends-and-predictions/}
\end{itemize}

%------------------------------------------------------------------------------
\newpage
\subsection{Graduate School Has Become a \$100,000 Panic Room for Students Who Cannot Face the Job Market}

\subsubsection*{Core Argument}
The surge in graduate school applications isn't intellectual ambition---it's a mass psychological retreat from economic reality. Students pay six-figure sums to delay confronting a hostile job market, often emerging with identical or worse prospects plus substantial debt.

\subsubsection*{Angles to Explore}
\begin{enumerate}[noitemsep]
    \item \textbf{The ROI Reckoning}: Nursing MSN vs. English MA---wildly different returns
    \item \textbf{The Adjunctification Warning}: 70\% of college instructors are contingent faculty
    \item \textbf{Professional Schools as Safer Bets?}: Law school oversupply, medical school barriers
    \item \textbf{The Alternative Path}: Professional certifications, bootcamps, strategic experience
\end{enumerate}

\subsubsection*{Key Data Points}
\begin{itemize}[noitemsep]
    \item Average graduate student debt: \$65,000+
    \item Time to recover grad school investment: 10-15 years for many humanities MAs
    \item Under 30\% of PhD graduates in tenure-track positions
    \item USC graduate program tuition ranges
\end{itemize}

\subsubsection*{Sources to Interview}
Dr. Leonard Cassuto (Fordham), USC Graduate School Dean, Graduate Career Consortium, financial advisors specializing in education debt, AAUP

\subsubsection*{Key URLs}
\begin{itemize}[noitemsep]
    \item \url{https://cgsnet.org/}
    \item \url{https://www.aaup.org/}
    \item \url{https://www.bls.gov/emp/chart-unemployment-earnings-education.htm}
    \item \url{https://gradcareerconsortium.org/}
\end{itemize}

%------------------------------------------------------------------------------
\newpage
\subsection{South Carolina Built an Economic Miracle That Has No Room for Its Own College Graduates}

\subsubsection*{Core Argument}
SC's economic transformation---BMW, Boeing, Volvo---is genuine achievement that has bypassed the state's university graduates almost entirely. Jobs driving GDP growth require technical certifications, not bachelor's degrees. SC has created two parallel economies, and USC trains students for the wrong one.

\subsubsection*{Angles to Explore}
\begin{enumerate}[noitemsep]
    \item \textbf{The Manufacturing Mismatch}: Midlands Tech graduates often out-earn USC Arts and Sciences graduates
    \item \textbf{The Brain Drain Bind}: Accept underemployment locally or migrate
    \item \textbf{Columbia's White-Collar Ceiling}: State government, healthcare, university---limited scale
    \item \textbf{What Adaptation Would Look Like}: Expand engineering, create hybrid credentials
\end{enumerate}

\subsubsection*{Key Data Points}
\begin{itemize}[noitemsep]
    \item SC GDP growth rate vs. bachelor's-required job creation
    \item USC graduation numbers by college vs. SC job openings
    \item Net migration of 25-34 year olds with degrees out of SC
    \item State appropriations per USC student vs. per technical college student
\end{itemize}

\subsubsection*{Sources to Interview}
SC Dept of Commerce economists, ReadySC leadership, USC Economic Development Office, Midlands Tech president, BMW/Boeing/Volvo HR, Columbia Chamber of Commerce

\subsubsection*{Key URLs}
\begin{itemize}[noitemsep]
    \item \url{https://www.sccommerce.com/}
    \item \url{https://www.bls.gov/regions/southeast/south_carolina.htm}
    \item \url{https://www.sctechsystem.edu/}
    \item \url{https://www.richmondfed.org/}
\end{itemize}

%------------------------------------------------------------------------------
\newpage
\subsection{The Internship Industrial Complex Is a Rigged Game That Punishes Students Who Cannot Afford to Play}

\subsubsection*{Core Argument}
The internship has transformed from supplementary experience into mandatory credential that systematically advantages wealthy students. When 48\% of jobs are in-person and underemployment hits 41.8\%, unpaid internships constitute a poll tax on career entry.

\subsubsection*{Angles to Explore}
\begin{enumerate}[noitemsep]
    \item \textbf{The Geography of Opportunity}: DC summer requires family subsidies most can't provide
    \item \textbf{The Credit Hour Scam}: Students pay tuition to work for free
    \item \textbf{The Experience Paradox}: Entry-level requires experience; experience requires being hired
    \item \textbf{What Equity Would Require}: Mandatory paid internships, university subsidies
\end{enumerate}

\subsubsection*{Key Data Points}
\begin{itemize}[noitemsep]
    \item $\sim$40\% of internships unpaid in certain sectors
    \item Cost of living in DC/NYC/Atlanta vs. average internship compensation
    \item Correlation between family income and internship completion
    \item 48\% of positions now fully in-person
\end{itemize}

\subsubsection*{Sources to Interview}
Ross Perlin (``Intern Nation''), USC Career Center, Financial Aid office, first-generation USC students, NACE, employers with paid-only models

\subsubsection*{Key URLs}
\begin{itemize}[noitemsep]
    \item \url{https://www.naceweb.org/}
    \item \url{https://www.dol.gov/agencies/whd/fact-sheets/71-flsa-internships}
    \item \url{https://cepr.net/}
    \item \url{https://payourinterns.org/}
\end{itemize}

%==============================================================================
% CATEGORY 7: POLITICS & 2026 MIDTERMS
%==============================================================================
\newpage
\section{Politics \& 2026 Midterms}

\textit{Context: Democrats hold 14-point lead on generic ballot. Trump at 38\% approval. Dems need only 3 House and 4 Senate seats. 228 executive orders in first year.}

%------------------------------------------------------------------------------
\newpage
\subsection{The Executive Order Presidency Has Made Congress Irrelevant}

\subsubsection*{Core Argument}
President Trump's 228 executive orders represent not merely aggressive governance but fundamental restructuring of American democracy that renders the legislative branch ceremonial. This expansion of executive power---regardless of party---should alarm students who will inherit a government where presidents rule by decree.

\subsubsection*{Angles to Explore}
\begin{enumerate}[noitemsep]
    \item \textbf{Constitutional Crisis in Plain Sight}: 228 orders vs. historical norms
    \item \textbf{Your Generation's Political Power}: Major decisions bypass representatives students can influence
    \item \textbf{Courts as New Legislature}: Policy made by president, unmade by courts
    \item \textbf{Both Parties Are Guilty}: Democrats expanded executive power under Obama and Biden
\end{enumerate}

\subsubsection*{Key Data Points}
\begin{itemize}[noitemsep]
    \item 228 executive orders in Trump's first year
    \item Congressional productivity: bills passed per session over 30 years
    \item Legal challenge success rates against executive orders
    \item Voter turnout in midterm vs. presidential elections
\end{itemize}

\subsubsection*{Sources to Interview}
Dr. William Howell (Chicago), Prof. Saikrishna Prakash (UVA Law), USC Political Science faculty, Protect Democracy, Federalist Society scholars

\subsubsection*{Key URLs}
\begin{itemize}[noitemsep]
    \item \url{https://www.brookings.edu/articles/tracking-regulatory-changes-in-the-biden-and-trump-administrations/}
    \item \url{https://www.presidency.ucsb.edu/statistics/data/executive-orders}
    \item \url{https://crsreports.congress.gov/}
    \item \url{https://protectdemocracy.org/}
\end{itemize}

%------------------------------------------------------------------------------
\newpage
\subsection{South Carolina Republicans Are Betting Their Future on Issues That Won't Save Them in November}

\subsubsection*{Core Argument}
While national Republicans face potential electoral catastrophe, SC GOP leaders are doubling down on DUI reform, tax cuts, and immigration. This reveals a party that has mistaken safe-state status for ideological validation, and may cost competitive House seats.

\subsubsection*{Angles to Explore}
\begin{enumerate}[noitemsep]
    \item \textbf{The Safe-State Delusion}: SC hasn't voted Democratic since 1976---breeding complacency
    \item \textbf{The Immigration Disconnect}: What do polls show SC voters actually prioritize?
    \item \textbf{The Nancy Mace Problem}: SC-01 competitive seat vs. state party's rightward drift
    \item \textbf{What a Blue Wave Would Mean}: Fewer committee assignments, less federal funding
\end{enumerate}

\subsubsection*{Key Data Points}
\begin{itemize}[noitemsep]
    \item Trump: 38\% approval; Dems: 14-point generic ballot lead
    \item SC-01 election margins: 2020, 2022, 2024
    \item Demographic shifts in Charleston County
    \item USC student voter registration in Richland County
\end{itemize}

\subsubsection*{Sources to Interview}
Dr. Scott Huffmon (Winthrop Poll), Gibbs Knotts (College of Charleston), USC College Democrats/Republicans, Cook Political Report analysts

\subsubsection*{Key URLs}
\begin{itemize}[noitemsep]
    \item \url{https://www.winthrop.edu/winthroppoll/}
    \item \url{https://www.cookpolitical.com/ratings/house-race-ratings}
    \item \url{https://www.scstatehouse.gov/}
    \item \url{https://projects.fivethirtyeight.com/polls/generic-ballot/}
\end{itemize}

%------------------------------------------------------------------------------
\newpage
\subsection{Gen Z Will Decide the Midterms But Most of Us Won't Bother Voting}

\subsubsection*{Core Argument}
Gen Z is now the largest potential voting bloc. The 2026 midterms represent this generation's first opportunity to reshape Congress. Yet historical patterns suggest most---including USC students---will stay home, effectively handing power to older generations whose interests diverge from ours.

\subsubsection*{Angles to Explore}
\begin{enumerate}[noitemsep]
    \item \textbf{The Math That Should Terrify}: If Gen Z voted at 65+ rates, we'd control virtually every competitive race
    \item \textbf{The Cynicism Trap}: ``It doesn't matter'' becomes self-fulfilling
    \item \textbf{The USC-Specific Problem}: Registration in Richland County vs. home; Greek life vs. off-campus
    \item \textbf{What Politicians Think About You}: Consultants advise deprioritizing youth outreach
\end{enumerate}

\subsubsection*{Key Data Points}
\begin{itemize}[noitemsep]
    \item Youth turnout never exceeded 36\% in midterms (2018 historic high)
    \item Democrats need 3 House seats, 4 Senate seats
    \item USC student population: 35,000+
    \item Student loan debt statistics
\end{itemize}

\subsubsection*{Sources to Interview}
CIRCLE (Tufts), John Della Volpe (Harvard), Campus Vote Project, USC Student Government, Richland County Elections Office, campaign operatives (anonymously)

\subsubsection*{Key URLs}
\begin{itemize}[noitemsep]
    \item \url{https://circle.tufts.edu/}
    \item \url{https://www.campusvoteproject.org/}
    \item \url{https://www.scvotes.gov/}
    \item \url{https://iop.harvard.edu/youth-poll}
\end{itemize}

%------------------------------------------------------------------------------
\newpage
\subsection{The Republican Retirement Wave Tells You Everything About Where This Party Is Headed}

\subsubsection*{Core Argument}
The exodus of senior Republican figures isn't typical turnover---it's mass evacuation from a party establishment conservatives no longer recognize. These retirements will accelerate the GOP's transformation into permanently Trumpist institution, with consequences extending decades.

\subsubsection*{Angles to Explore}
\begin{enumerate}[noitemsep]
    \item \textbf{Names That Matter}: Profile significant retirements---age/term limits vs. fleeing party direction
    \item \textbf{Institutional Knowledge Loss}: Expertise in policy, procedure, compromise walking out
    \item \textbf{Who Replaces Them}: Primary dynamics in safe Republican seats
    \item \textbf{The Long Game}: 2028 implications and beyond
\end{enumerate}

\subsubsection*{Key Data Points}
\begin{itemize}[noitemsep]
    \item Number of Republican vs. Democratic retirements
    \item Historical comparison to 1994, 2010, 2018 waves
    \item Ideological scores (DW-NOMINATE) of retiring members
    \item Trump endorsement success rate in primaries
\end{itemize}

\subsubsection*{Sources to Interview}
Retiring Republican members/staff, former Republican officials who left party, Tim Miller (The Bulwark), conservative think tank scholars, USC political science faculty

\subsubsection*{Key URLs}
\begin{itemize}[noitemsep]
    \item \url{https://www.cookpolitical.com/}
    \item \url{https://ballotpedia.org/}
    \item \url{https://voteview.com/}
    \item \url{https://www.thebulwark.com/}
\end{itemize}

%==============================================================================
% CATEGORY 8: HEALTHCARE & PUBLIC HEALTH
%==============================================================================
\newpage
\section{Healthcare \& Public Health}

\textit{Context: ACA subsidies expired---20M seeing 114\% premium increases. 4.8M projected to drop coverage. US could lose measles-free status. 2,012 measles cases in 2025. CDC weakened by layoffs.}

%------------------------------------------------------------------------------
\newpage
\subsection{The ACA Cliff Will Push a Generation of College Students Into Medical Debt}

\subsubsection*{Core Argument}
The expiration of ACA subsidies represents the largest deliberate transfer of healthcare costs onto young Americans in a generation. With 4.8 million projected to lose coverage and 114\% premium increases, college students will choose between health insurance and basic necessities.

\subsubsection*{Angles to Explore}
\begin{enumerate}[noitemsep]
    \item \textbf{The ``Aging Out'' Crisis}: Students turning 26 lose parental coverage into hostile marketplace
    \item \textbf{USC's Student Health Insurance}: Costs vs. marketplace; what if students can't afford it?
    \item \textbf{Hidden Costs of Being Uninsured at 22}: Delayed care creates compounding problems
    \item \textbf{SC's Medicaid Gap}: Non-expansion state---students fall into coverage gap
\end{enumerate}

\subsubsection*{Key Data Points}
\begin{itemize}[noitemsep]
    \item 20 million seeing 114\% premium increases
    \item 4.8 million projected to drop coverage
    \item SC uninsured rate: 11.2\% (above 8.6\% national average)
    \item SC is one of 10 states without Medicaid expansion
\end{itemize}

\subsubsection*{Sources to Interview}
Dr. Ana Langer (Harvard), USC Student Health Services Director, SC Appleseed Legal Justice Center, Young Invincibles, Dr. Benjamin Sommers (Harvard)

\subsubsection*{Key URLs}
\begin{itemize}[noitemsep]
    \item \url{https://www.kff.org/health-reform/}
    \item \url{https://www.sa.sc.edu/shs/}
    \item \url{https://www.commonwealthfund.org/}
    \item \url{https://younginvincibles.org/}
\end{itemize}

%------------------------------------------------------------------------------
\newpage
\subsection{The Return of Measles to America Is a Policy Choice, Not a Medical Mystery}

\subsubsection*{Core Argument}
2,012 measles cases in 2025---most since early 1990s---is not failure of science or inevitable consequence of vaccine hesitancy. It's the direct result of policy decisions: CDC cuts, legitimization of anti-vaccine rhetoric, weakening of public health infrastructure.

\subsubsection*{Angles to Explore}
\begin{enumerate}[noitemsep]
    \item \textbf{USC's Vaccination Requirements Under Pressure}: SC allows religious/philosophical exemptions
    \item \textbf{The CDC Gutting}: Layoffs decimated surveillance, outbreak response
    \item \textbf{The ``Vaccine Schedule Overhaul'' Threat}: Political rather than scientific considerations
    \item \textbf{International Student Dimension}: Robust requirements ensure everyone protected
\end{enumerate}

\subsubsection*{Key Data Points}
\begin{itemize}[noitemsep]
    \item 2,012 measles cases in 2025---highest since early 1990s
    \item US at risk of losing measles elimination status (achieved 2000)
    \item Herd immunity requires 95\% vaccination; SC kindergarten: $\sim$92\%
    \item Measles: 90\% infection rate among unvaccinated exposed
\end{itemize}

\subsubsection*{Sources to Interview}
Dr. Peter Hotez (Baylor), Dr. Paul Offit (CHOP), USC Student Health immunization coordinator, SC DHEC epidemiology, American College Health Association

\subsubsection*{Key URLs}
\begin{itemize}[noitemsep]
    \item \url{https://www.cdc.gov/measles/}
    \item \url{https://scdhec.gov/health/vaccines}
    \item \url{https://www.sa.sc.edu/shs/immunizations/}
    \item \url{https://www.acha.org/}
\end{itemize}

%------------------------------------------------------------------------------
\newpage
\subsection{AI in Healthcare Will Not Save Us---It Will Deepen the Divide}

\subsubsection*{Core Argument}
80\% of health executives expect AI to ``deliver value''---but value for whom? AI tools will first deploy in wealthy health systems, widening the gap between haves and have-nots. For students at a public university in a non-expansion state, AI promises a two-tiered system.

\subsubsection*{Angles to Explore}
\begin{enumerate}[noitemsep]
    \item \textbf{The Bias Problem}: AI trained on white, wealthy patient data performs poorly for Black, rural, low-income patients
    \item \textbf{Palmetto GBA and AI in SC}: How tools being deployed; what oversight exists
    \item \textbf{The Mental Health AI Mirage}: Chatbots as cheap substitute for actual counseling
    \item \textbf{Data Privacy}: Who owns student health app data?
\end{enumerate}

\subsubsection*{Key Data Points}
\begin{itemize}[noitemsep]
    \item 80\% of health executives expect AI to deliver value
    \item Optum algorithm: Black patients received lower risk scores than equally sick white patients
    \item SC: 60\% designated Mental Health Professional Shortage Area
    \item Only 40\% of AI health tools have peer-reviewed validation
\end{itemize}

\subsubsection*{Sources to Interview}
Dr. Ziad Obermeyer (UC Berkeley), Dr. Ruha Benjamin (Princeton), USC Arnold School faculty, Electronic Frontier Foundation, SC Hospital Association

\subsubsection*{Key URLs}
\begin{itemize}[noitemsep]
    \item \url{https://www.science.org/doi/10.1126/science.aax2342}
    \item \url{https://www.eff.org/issues/medical-privacy}
    \item \url{https://www.nature.com/nm/}
    \item \url{https://ainowinstitute.org/}
\end{itemize}

%------------------------------------------------------------------------------
\newpage
\subsection{South Carolina's Healthcare Desert Is Not Geography---It Is Politics}

\subsubsection*{Core Argument}
When rural hospitals close and students from small-town SC arrive at USC without primary care, we call it a ``healthcare desert''---as if natural. It is not. SC's healthcare landscape is direct result of political choices: Medicaid non-expansion, certificate-of-need laws, legislative ideology over constituent health.

\subsubsection*{Angles to Explore}
\begin{enumerate}[noitemsep]
    \item \textbf{Rural Hospital Death Spiral}: Map where USC students come from; overlay hospital closures
    \item \textbf{Medicaid Expansion}: Calculate how many USC students/families would gain coverage
    \item \textbf{Primary Care Pipeline Problem}: Physician shortage tied to state policy
    \item \textbf{Student Health as Bandaid}: Campus health never designed for students with years of deferred care
\end{enumerate}

\subsubsection*{Key Data Points}
\begin{itemize}[noitemsep]
    \item SC rural hospital closures (Chartis Center data)
    \item Medicaid expansion would cover 200,000+ additional SC residents
    \item SC primary care physician shortage: $\sim$400+ below need
    \item Majority of SC counties designated Health Professional Shortage Areas
\end{itemize}

\subsubsection*{Sources to Interview}
SC Appleseed, Dr. Graham Adams (USC School of Medicine), SC Rural Health Association, state legislators, USC students from rural counties, Chartis Center for Rural Health

\subsubsection*{Key URLs}
\begin{itemize}[noitemsep]
    \item \url{https://www.kff.org/medicaid/}
    \item \url{https://www.chartis.com/services/community-health/the-chartis-center-for-rural-health}
    \item \url{https://www.scstatehouse.gov/}
    \item \url{https://www.scha.org/}
\end{itemize}

%==============================================================================
% CATEGORY 9: GLOBAL AFFAIRS & AMERICAN ISOLATIONISM
%==============================================================================
\newpage
\section{Global Affairs \& American Isolationism}

\textit{Context: US withdrew from 66 international organizations. Venezuela operation. Greenland threats. Ukraine losing ground. US considered ``principal source of global risk.'' International student interest dropped 47\%.}

%------------------------------------------------------------------------------
\newpage
\subsection{America's Retreat from the World Will Cost This Generation More Than Any Other}

\subsubsection*{Core Argument}
Withdrawal from 66 international organizations represents not policy shift but generational betrayal. Young Americans entering the workforce inherit a world where American influence has been deliberately dismantled, American credentials carry less weight, and international systems that enabled prosperity have been undermined.

\subsubsection*{Angles to Explore}
\begin{enumerate}[noitemsep]
    \item \textbf{The Credential Devaluation Crisis}: American degrees derived value from American-led institutions
    \item \textbf{The Soft Power Inheritance}: Previous generations inherited American influence; this generation inherits aftermath
    \item \textbf{The Knowledge Economy Collapse}: Scientific collaboration reduced; brain drain begun
    \item \textbf{The Democracy Deficit}: Power vacuums filled by authoritarian regimes
\end{enumerate}

\subsubsection*{Key Data Points}
\begin{itemize}[noitemsep]
    \item 47\% drop in international student enrollment interest
    \item 66 international organizations withdrawn from
    \item Economic cost of lost soft power
    \item USC study abroad participation rates
\end{itemize}

\subsubsection*{Sources to Interview}
Dr. Joseph Nye (Harvard), USC Arnold School faculty, USC Career Center, Foreign Policy analysts, Council on Foreign Relations scholars

\subsubsection*{Key URLs}
\begin{itemize}[noitemsep]
    \item \url{https://opendoorsdata.org/}
    \item \url{https://www.cfr.org/}
    \item \url{https://softpower30.com/}
    \item \url{https://www.foreignaffairs.com/}
\end{itemize}

%------------------------------------------------------------------------------
\newpage
\subsection{The 350,000 Dead Are Not an Abstraction and American Students Must Reckon With What Their Government Has Done}

\subsubsection*{Core Argument}
The estimated 350,000+ deaths attributed to American aid cuts represent moral catastrophe demanding honest confrontation. USC students---future doctors, public health officials, diplomats---must understand these are consequences of deliberate choices, not statistics.

\subsubsection*{Angles to Explore}
\begin{enumerate}[noitemsep]
    \item \textbf{Moral Mathematics}: PEPFAR cuts and HIV treatment interruption; WHO withdrawal and surveillance gaps
    \item \textbf{The Complicity Question}: What does citizenship mean when your government causes mass death?
    \item \textbf{The Memory Hole Problem}: Minimal American media coverage---why?
    \item \textbf{Career Implications of Conscience}: How to navigate institutions responsible for such outcomes?
\end{enumerate}

\subsubsection*{Key Data Points}
\begin{itemize}[noitemsep]
    \item 350,000+ estimated deaths from aid cuts
    \item WHO budget shortfall and disease surveillance gaps
    \item PEPFAR funding: peak vs. current
    \item Sudan civil war death toll and US response reduction
\end{itemize}

\subsubsection*{Sources to Interview}
Dr. Paul Farmer's writings (posthumously), Physicians for Human Rights, USC School of Medicine faculty, former USAID officials who resigned, Partners in Health

\subsubsection*{Key URLs}
\begin{itemize}[noitemsep]
    \item \url{https://www.kff.org/global-health-policy/}
    \item \url{https://phr.org/}
    \item \url{https://www.thelancet.com/journals/langlo/home}
    \item \url{https://www.who.int/about/funding}
\end{itemize}

%------------------------------------------------------------------------------
\newpage
\subsection{International Students Are Fleeing American Universities and Taking the Future With Them}

\subsubsection*{Core Argument}
The 47\% collapse in international student enrollment interest is not an admissions problem---it's a leading indicator of American decline. International students represent future business partners, diplomatic contacts, cultural ambassadors. Their departure to Canadian, British, Chinese universities represents generational transfer of influence.

\subsubsection*{Angles to Explore}
\begin{enumerate}[noitemsep]
    \item \textbf{Human Stories}: Profile specific international students at USC navigating visa uncertainty
    \item \textbf{Network Effect Loss}: Relationships that structure future international business disappearing
    \item \textbf{Campus Culture Transformation}: Diverse perspectives diminishing; American students receive provincial education
    \item \textbf{Competitor Advantage}: Canada, Australia, UK, China actively recruiting fleeing students
\end{enumerate}

\subsubsection*{Key Data Points}
\begin{itemize}[noitemsep]
    \item 47\% decline in international student interest
    \item USC international enrollment: historical trends, projected decline
    \item Economic contribution of international students
    \item Visa denial rates and processing time increases
\end{itemize}

\subsubsection*{Sources to Interview}
USC International Student Services, current international students, students who transferred out, admissions officers at competitor institutions, NAFSA

\subsubsection*{Key URLs}
\begin{itemize}[noitemsep]
    \item \url{https://www.nafsa.org/policy-and-advocacy/policy-resources/nafsa-international-student-economic-value-tool}
    \item \url{https://opendoorsdata.org/}
    \item \url{https://www.timeshighereducation.com/}
    \item \url{https://www.migrationpolicy.org/}
\end{itemize}

%------------------------------------------------------------------------------
\newpage
\subsection{From Greenland Threats to Gaza Silence We Are Watching American Moral Authority Die in Real Time}

\subsubsection*{Core Argument}
The combination of territorial threats against allies (Greenland), complicity in humanitarian catastrophe (Gaza), abandonment of partners (Ukraine), and institutional withdrawal represents the death of American moral authority---the intangible foundation of American global leadership.

\subsubsection*{Angles to Explore}
\begin{enumerate}[noitemsep]
    \item \textbf{The Greenland Episode}: Coercion against NATO ally reveals willingness that undermines alliances
    \item \textbf{The Gaza Complicity}: American diplomatic cover produces global perception shifts
    \item \textbf{The Ukraine Abandonment}: Demonstrates to every ally that commitments are conditional
    \item \textbf{The Moral Authority Economy}: Once lost, requires generations to rebuild
\end{enumerate}

\subsubsection*{Key Data Points}
\begin{itemize}[noitemsep]
    \item Global opinion polling on American trustworthiness (Pew, Gallup)
    \item Gaza casualty figures and humanitarian indicators
    \item Ukraine territorial losses correlated with American support reduction
    \item ``Principal source of global risk'' polling
\end{itemize}

\subsubsection*{Sources to Interview}
European foreign policy analysts, USC international relations faculty, Palestinian and Israeli students at USC, Ukrainian students, think tank scholars (Brookings, Carnegie, CSIS)

\subsubsection*{Key URLs}
\begin{itemize}[noitemsep]
    \item \url{https://www.pewresearch.org/global/}
    \item \url{https://www.gallup.com/analytics/318875/global-research.aspx}
    \item \url{https://ecfr.eu/}
    \item \url{https://carnegieendowment.org/}
\end{itemize}

%==============================================================================
% CATEGORY 10: SOUTH CAROLINA & USC LOCAL ISSUES
%==============================================================================
\newpage
\section{South Carolina \& USC Local Issues}

\textit{Context: SC has highest DUI fatalities per capita. Population growing 90,000/year. Housing affordability concerns. Women's basketball \#3 nationally. NIL transforming athletics. No tuition increase for 7th year.}

%------------------------------------------------------------------------------
\newpage
\subsection{The University Froze Tuition for Seven Years While Students Still Cannot Afford to Live Here}

\subsubsection*{Core Argument}
USC's celebrated seven-year tuition freeze is a masterclass in misdirection---a PR victory obscuring the real affordability crisis. While administrators take victory laps, actual cost of attending USC has exploded due to skyrocketing Columbia rents, mandatory fees, and auxiliary costs.

\subsubsection*{Angles to Explore}
\begin{enumerate}[noitemsep]
    \item \textbf{The Hidden Cost Explosion}: Mandatory fees, housing up even as tuition flat
    \item \textbf{Columbia's Rent Crisis}: Vista, Five Points, Olympia transformed into premium markets
    \item \textbf{Political Theater of Tuition Freezes}: Why popular despite limited impact
    \item \textbf{Real Affordability}: Rent guarantee programs, affordable housing partnerships
\end{enumerate}

\subsubsection*{Key Data Points}
\begin{itemize}[noitemsep]
    \item USC tuition frozen at $\sim$\$12,688 since 2018
    \item Columbia rent increased 25-35\% since 2018
    \item USC cost of attendance exceeds \$28,000 for in-state
    \item Mandatory fees add \$800+ annually
    \item SC population growth: 90,000/year
\end{itemize}

\subsubsection*{Sources to Interview}
Dr. Christy Friend (Provost Office), Columbia City Council, Dr. Bobby Donaldson (USC History), USC Student Government, SC Commission on Higher Education

\subsubsection*{Key URLs}
\begin{itemize}[noitemsep]
    \item \url{https://sc.edu/about/offices_and_divisions/financial_aid/cost_of_attendance/}
    \item \url{https://www.che.sc.gov/}
    \item \url{https://www.apartmentlist.com/renter-life/average-rent-in-columbia-sc}
    \item \url{https://www.sc.edu/about/offices_and_divisions/secretary/toolbox/trustees/}
\end{itemize}

%------------------------------------------------------------------------------
\newpage
\subsection{South Carolina's Social Media Age Verification Law Will Fail Students It Claims to Protect}

\subsubsection*{Core Argument}
SC's new social media age verification law is constitutional theater that will accomplish nothing except normalizing government surveillance. The law cannot be enforced, will be struck down on First Amendment grounds, and distracts from evidence-based mental health approaches.

\subsubsection*{Angles to Explore}
\begin{enumerate}[noitemsep]
    \item \textbf{The Constitutional Dead End}: Arkansas, Ohio, Texas laws blocked on First Amendment grounds
    \item \textbf{The Privacy Paradox}: Age verification requires invasive data collection
    \item \textbf{What Actually Affects Student Mental Health}: More counseling staff, shorter wait times
    \item \textbf{A Generation's Perspective}: Students as active participants, not passive victims
\end{enumerate}

\subsubsection*{Key Data Points}
\begin{itemize}[noitemsep]
    \item Federal court injunction rates for similar state laws
    \item USC Counseling wait times and staffing ratios
    \item CDC Youth Risk Behavior Survey data
    \item Cost to platforms/state for compliance
\end{itemize}

\subsubsection*{Sources to Interview}
Dr. Amanda Fairchild (USC Psychology), Dr. Scott Babcock (Student Health), Electronic Frontier Foundation, USC Law faculty, SC legislators, NetChoice

\subsubsection*{Key URLs}
\begin{itemize}[noitemsep]
    \item \url{https://www.scstatehouse.gov/}
    \item \url{https://www.eff.org/issues/age-verification}
    \item \url{https://netchoice.org/}
    \item \url{https://www.cdc.gov/yrbs/}
\end{itemize}

%------------------------------------------------------------------------------
\newpage
\subsection{Dawn Staley's Program Proves NIL Can Build Dynasties Without Destroying College Sports}

\subsubsection*{Core Argument}
The hand-wringing about NIL destroying college athletics ignores evidence at Colonial Life Arena: Dawn Staley's \#3-ranked women's basketball program demonstrates that NIL, properly integrated, strengthens rather than corrupts college sports.

\subsubsection*{Angles to Explore}
\begin{enumerate}[noitemsep]
    \item \textbf{The Women's Basketball NIL Model}: How Staley's program integrated NIL across roster
    \item \textbf{Why Women's Sports NIL Is Different}: Local partnerships, community engagement vs. bidding wars
    \item \textbf{Culture as Competitive Advantage}: Sustained excellence, player development, low transfer rates
    \item \textbf{The USC Athletics Disparity}: Contrast with football's more turbulent experience
\end{enumerate}

\subsubsection*{Key Data Points}
\begin{itemize}[noitemsep]
    \item USC women's basketball: \#3 ranking
    \item Player retention rates vs. national averages
    \item NIL valuations for top USC women's players (On3)
    \item Colonial Life Arena attendance growth
    \item Dawn Staley: 3 National Championships, 7 Final Fours in last 10 tournaments
\end{itemize}

\subsubsection*{Sources to Interview}
Dawn Staley, USC Athletic Director, current USC women's basketball players, Garnet Trust representatives, Dr. Victoria Jackson (ASU sports historian)

\subsubsection*{Key URLs}
\begin{itemize}[noitemsep]
    \item \url{https://gamecocksonline.com/sports/womens-basketball}
    \item \url{https://www.on3.com/nil/}
    \item \url{https://www.ncaa.org/sports/2021/6/28/name-image-likeness.aspx}
    \item \url{https://theathletic.com/college-basketball/}
\end{itemize}

%------------------------------------------------------------------------------
\newpage
\subsection{The Fight Over Campus Speakers at USC Reveals That Nobody Actually Believes in Free Speech}

\subsubsection*{Core Argument}
Recurring controversies over speaker events expose fundamental dishonesty across the political spectrum: everyone claims to support free speech for their side. USC should model genuine viewpoint diversity by requiring intellectual consistency.

\subsubsection*{Angles to Explore}
\begin{enumerate}[noitemsep]
    \item \textbf{Anatomy of a Speaker Controversy}: Examine specific 2025 USC events
    \item \textbf{The Selective Outrage Problem}: Document how same individuals apply different standards
    \item \textbf{The Security Cost Weapon}: Controversial speech becomes expensive speech
    \item \textbf{Genuine Viewpoint Diversity}: Either equal support for all legal speech, or openly acknowledge value judgments
\end{enumerate}

\subsubsection*{Key Data Points}
\begin{itemize}[noitemsep]
    \item Recent speaker events at USC and outcomes
    \item USC speaker policies
    \item FIRE ratings for USC
    \item Security cost figures for controversial events
\end{itemize}

\subsubsection*{Sources to Interview}
USC student organization leaders, Suzanne Nossel (PEN America), FIRE representatives, USC administration, students who protested speakers, USC Law faculty

\subsubsection*{Key URLs}
\begin{itemize}[noitemsep]
    \item \url{https://www.thefire.org/research-learn/campus-disinvitation-database}
    \item \url{https://pen.org/issue/free-expression-and-education/}
    \item \url{https://knightfoundation.org/topics/free-expression/}
    \item \url{https://www.dailygamecock.com/}
\end{itemize}

%==============================================================================
% APPENDIX: WRITING GUIDELINES
%==============================================================================
\newpage
\section*{Appendix: Writing Guidelines for Opinion Pieces}

\subsection*{The Daily Gamecock Voice}
\begin{itemize}[noitemsep]
    \item \textbf{Argumentative, not reportorial}: Pass judgment immediately rather than explaining neutrally
    \item \textbf{Conversational but rigorous}: Vivid imagery paired with hard data
    \item \textbf{Personal but connected}: First-person experience as entry point to systemic critique
    \item \textbf{Urgent but measured}: Conveys seriousness without false urgency
\end{itemize}

\subsection*{Headline Rules}
\begin{itemize}[noitemsep]
    \item No colons in headlines---use declarative sentences that make arguments
    \item No false urgency (avoid ``Just,'' ``Now,'' ``Recently'')
    \item Make blame explicit with strong verbs (``sabotage,'' ``exploit,'' ``betray'')
\end{itemize}

\subsection*{Structure (1,200--1,400 words)}
\begin{itemize}[noitemsep]
    \item Introduction with student anecdote: 150--200 words
    \item Each angle: 250--300 words
    \item Conclusion with synthesis: 150--200 words
\end{itemize}

\subsection*{Evergreen Approach}
For a monthly print magazine, timely subjects should be analyzed through evergreen lenses:
\begin{itemize}[noitemsep]
    \item Frame around structural analysis rather than news events
    \item Include historical context and patterns
    \item Provide frameworks readers can apply to future situations
    \item End with calls to action or reflection questions
\end{itemize}

\subsection*{Sourcing Requirements}
\begin{itemize}[noitemsep]
    \item 3--5 USC student voices (diverse backgrounds)
    \item At least one national expert and one state/local source per piece
    \item Request official data from USC administration
    \item Document all attempts to reach sources
\end{itemize}

\end{document}
